\documentclass[10pt,a4paper]{amsart}
\usepackage[margin=19mm]{geometry}
\usepackage{amsmath,amssymb,mathtools}
\usepackage[scr=rsfs,cal=euler]{mathalfa}
\usepackage{xfrac}
\usepackage[unicode,hidelinks,bookmarks=false]{hyperref}
\newcommand{\ie}{i.e.\ }
\newcommand{\cf}{cf.\ }
\newcommand{\resp}{respectively\ }
\newcommand{\Subsection}{\subsection}
\providecommand{\nobreakdash}{\nobreak\hbox{-}}
\setlength{\emergencystretch}{2em}
\multlinegap0pt
\usepackage{amssymb}
\usepackage{algorithm}\usepackage{algpseudocode}
\usepackage{float}
\usepackage{xcolor}
\newtheorem{theorem}{Theorem}
\title{SIPF-PIC: An efficient Stochastic Interacting Particle-Field method for 3D fully parabolic Chemotaxis Systems}
\author{Jingyuan Hu, Zhongjian Wang, Jack Xin, Zhiwen Zhang}
\date{Source: arXiv:2512.03452v4 (CC BY 4.0)}
\begin{document}
\maketitle

\noindent\emph{Source and licence.} SIPF-PIC: An efficient Stochastic Interacting Particle-Field method for 3D fully parabolic Chemotaxis Systems. Version arXiv:2512.03452v4 (CC BY 4.0), distributed by arXiv under the Creative Commons Attribution 4.0 International licence, http://creativecommons.org/licenses/by/4.0/, as recorded in the arXiv licence metadata for this exact version and shown on the abstract page https://arxiv.org/abs/2512.03452v4. Full licence notice: \texttt{licenses/arxiv-2512.03452-CC-BY-4.0.txt}.\par\smallskip

\noindent\textit{Editorial scope.} Counted unit: the article's theorem theorem:positive\_diffusion\_kernel (source0.tex lines 690--708). The article's complete original proof of it is delivered with it (source0.tex lines 710--817, one \texttt{proof} environment of 108 lines). No mathematics is written, repaired or re-derived: the statement and the proof are the article's own lines, in the article's own notation, and no retained line is substituted or re-flowed.  No further source-local context is needed: every cross-reference of every retained line resolves inside the two delivered spans.  Nothing was omitted from any retained span, so the package carries no omission record.  No retained line contains a citation, so the wrapper carries no bibliography and no bibliography entry.  No retained line contains a figure environment, an image-inclusion command or drawing code, so no image is delivered and the package declares no asset dependency.  Editorial formatting only, in addition: an AMS \texttt{amsart} preamble that restates the article's own declarations and macros used by the retained lines, namely the environments \texttt{theorem} on separate counters, together with the standard packages those lines need.  The retained environments print in this document as Theorem 1 in order of appearance, whereas the article prints the same items with the numbering of its own full document.  The complete native source of the article as distributed in the arXiv e-print for this exact version is bundled unchanged as \texttt{sources/190264/source0.tex}.
\begin{theorem}
\label{theorem:positive_diffusion_kernel}
There exists a kernel
$\widetilde{G}_{\tau}\in\mathbb{R}^{U_H}$ whose components are strictly positive and whose discrete Fourier coefficients
$\widetilde{\kappa}_{\tau}=\mathcal{F}_H\widetilde{G}_{\tau}$ satisfy
\begin{equation}
\left|
\widetilde{\kappa}_{\tau,\mathbf q}
-
\hat{\kappa}_{\tau,\mathbf q}
\right|
\leq
4\exp\left(
-\frac{\pi^2\tau H^2}{\epsilon L^2}
\right),
\qquad \mathbf q\in U_H,
\end{equation}
provided that $\tau H^2L^{-2}$ is sufficiently large, i.e., $\tau\gtrsim (L/H)^{2-\sigma}$ for some $\sigma>0$.
\end{theorem}

\begin{proof}
For each $\mathbf q\in U_H$, define the aliased Fourier coefficients by
\begin{equation}
\widetilde{\kappa}_{\tau,\mathbf q}
=
\sum_{\mathbf n\in\mathbb Z^3}
\hat{\kappa}_{\tau,\mathbf q+\mathbf nH}
=
\sum_{\mathbf n\in\mathbb Z^3}
\exp\left(
-\frac{\tau}{\epsilon}
\left(
\left|
\frac{2\pi}{L}
(\mathbf q+\mathbf nH)
\right|^2
+k^2
\right)
\right),
\end{equation}
and let
$\widetilde{G}_{\tau}
=
\mathcal{F}_H^{-1}\widetilde{\kappa}_{\tau}$.
By the Poisson summation formula, the inverse discrete Fourier transform of
$\widetilde{\kappa}_{\tau}$ is obtained by sampling the periodization of the
corresponding Gaussian heat kernel. More precisely, for every grid point
$\mathbf x_{\mathbf p}=\mathbf pL/H$,
\begin{equation}
\widetilde{G}_{\tau}
\left(
\mathbf p\frac{L}{H}
\right)
=
\sum_{\mathbf n\in\mathbb Z^3}
G_{\tau}
\left(
\mathbf p\frac{L}{H}+\mathbf nL
\right),
\end{equation}
up to the normalization convention of $\mathcal{F}_H$, where
\begin{equation}
G_{\tau}(\mathbf x)
=
\exp\left(-\frac{k^2\tau}{\epsilon}\right)
\left(\frac{\epsilon}{4\pi\tau}\right)^{\frac{3}{2}}
\exp\left(
-\frac{\epsilon|\mathbf x|^2}{4\tau}
\right).
\end{equation}
Since every term in this sum is strictly positive, all components of
$\widetilde{G}_{\tau}$ are strictly positive.

It remains to estimate the difference between the aliased coefficients and
their principal components. For every $\mathbf q\in U_H$,
\begin{equation}
\begin{aligned}
\left|
\widetilde{\kappa}_{\tau,\mathbf q}
-
\hat{\kappa}_{\tau,\mathbf q}
\right|
&=
\sum_{\mathbf n\in\mathbb Z^3\setminus\{\mathbf0\}}
\exp\left(
-\frac{\tau}{\epsilon}
\left(
\left|
\frac{2\pi}{L}
(\mathbf q+\mathbf nH)
\right|^2
+k^2
\right)
\right)
\\
&\leq
\sum_{\mathbf n\in\mathbb Z^3\setminus\{\mathbf0\}}
\exp\left(
-\frac{\pi^2\tau H^2}{\epsilon L^2}
|\mathbf n|^2
\right).
\end{aligned}
\end{equation}
Writing $a=\pi^2\tau H^2/(\epsilon L^2)$, the remaining Gaussian lattice
tail satisfies
\begin{equation}
\sum_{\mathbf n\in\mathbb Z^3\setminus\{\mathbf0\}}
e^{-a|\mathbf n|^2}
\leq 4e^{-a}
\end{equation}
whenever $a$ is sufficiently large. This proves the stated estimate.

Finally, if
$\tau\gtrsim(L/H)^{2-\sigma}$, then
$\tau H^2/L^2\gtrsim H^\sigma$, and hence
\begin{equation}
\left|
\widetilde{\kappa}_{\tau,\mathbf q}
-
\hat{\kappa}_{\tau,\mathbf q}
\right|
\leq
4e^{-CH^\sigma}
\end{equation}
for some constant $C>0$. Since $e^{-CH^\sigma} \lesssim H^{-r}$ for every
$r>0$, this perturbation is asymptotically smaller than any polynomial
discretization error. Consequently, from an equivalent perspective, the explicitly computable kernel $\widehat{G}_{\tau}$ combined with the positivity-preserving interpolation operator $I$ introduces only an exponentially small violation of positivity in the computed concentration $\hat{c}$.
\end{proof}

\end{document}
